%% kapa: exam=1 round=34 year=2023 period=1 form=A subject=real_estate subject_ko=부동산학원론 seq=19 number=99
%% source: Q-Net 감정평가사 기출문제 (cst00309) artlSeq=5211789 file=raw/1008/5211789/2023년도 제34회 감정평가사 1차시험 1교시.pdf
%% conversion: kordoc 4.24.51 -> tools/gen1.py (auto); review=proofread-sonnet; flags=none
\begin{question}{99}
화폐의 시간가치에 관한 설명으로 옳지 않은 것은?
\begin{choices}
  \item 인플레이션, 화폐의 시차선호, 미래의 불확실성은 화폐의 시간가치를 발생시키는 요인이다.
  \item 감채기금이란 일정기간 후에 일정금액을 만들기 위해 매 기간 납입해야 할 금액을 말한다.
  \item 연금의 미래가치란 매 기간 마다 일정금액을 불입해 나갈 때, 미래 일정시점에서의 불입금액 총액의 가치를 말한다.
  \item 현재가치에 대한 미래가치를 산출하기 위하여 사용하는 이율을 이자율이라 하고, 미래가치에 대한 현재가치를 산출하기 위하여 사용하는 이율을 할인율이라 한다.
  \item 부동산 경기가 침체하는 시기에 상업용 부동산의 수익이 일정함에도 불구하고 부동산 가격이 떨어지는 것은 할인율이 낮아지기 때문이다.
\end{choices}
\end{question}
\answer{5}
